% =====================================================================
\chapter*{Instrucciones}
\addcontentsline{toc}{chapter}{Instrucciones}
\markboth{Instrucciones}{}
% =====================================================================

Este examen cierra la fase~III del curso, dedicada al Transformer. Cubre los módulos~12 (tokenización y modelos secuenciales), 13 (atención), 14 (el Transformer decoder-only) y 15 (entrenamiento de un mini-GPT). Debe aprobarse antes de comenzar la fase~IV.

\section*{Condiciones}
\begin{center}
\begin{tabular}{@{}ll@{}}
\toprule
Duración & 2~h 30~min: unos 20~min para la parte~A y 2~h para la parte~B \\
Material permitido & calculadora científica y una hoja de fórmulas manuscrita (una cara) \\
Puntuación total & 100 puntos: parte~A, 20; parte~B, 80 \\
Aprobado & 70 puntos (70\,\%) \\
Módulos & M12, M13, M14 y M15 \\
\bottomrule
\end{tabular}
\end{center}

\section*{Estructura}
\begin{itemize}[leftmargin=1.4em,label=\textendash]
  \item \textbf{Parte A. Cuestiones} (20 puntos). Veinte preguntas de opción múltiple, de 1~punto cada una, con una sola respuesta correcta entre las cuatro propuestas. Las respuestas en blanco o incorrectas valen~0; no se penalizan los errores.
  \item \textbf{Parte B. Problemas} (80 puntos). Seis ejercicios de desarrollo, con la puntuación de cada apartado indicada entre paréntesis. Debe justificarse cada paso: un resultado sin razonamiento vale como máximo la mitad de su puntuación.
  \item \textbf{Proyecto de fase}. Se entrega aparte y se evalúa con la rúbrica del capítulo \hyperref[chap:proyecto]{«Proyecto de fase»}; no suma puntos al examen.
\end{itemize}

\section*{Convenios}
\begin{itemize}[leftmargin=1.4em,label=\textendash]
  \item Se redondea a tres decimales salvo que se indique otra cosa. Pueden usarse $e\approx 2{,}718$, $e^{2}\approx 7{,}389$ y $\ln 2\approx 0{,}693$.
  \item Las matrices de pesos multiplican por la derecha ($Y=XW$) y no hay sesgos salvo que se diga lo contrario.
  \item $1~\mathrm{GB}=10^{9}$~bytes. Una multiplicación seguida de una suma cuenta como 2~FLOPs.
  \item Los modelos con nombre inventado, como «Lince-1B», son ficticios. Los datos de GPT-2 y de LLaMA-2 son los publicados por sus autores.
  \item Las referencias del tipo «M13, §2.3» remiten al módulo y a la sección de los apuntes del curso.
\end{itemize}

\section*{Sobre las fuentes}
Los enunciados son de redacción propia. Los ejercicios marcados con \emph{Fuente} están adaptados de tareas y ejemplos de cursos universitarios publicados: la tarea~1 de Stanford CS336~\cite{cs336}, la tarea~3 de Stanford CS224N~\cite{cs224n} y los ejercicios del libro de Jurafsky y Martin~\cite{jm3}, texto de los cursos de procesamiento del lenguaje de Stanford. En todos los casos se han cambiado los datos para que el ejercicio pueda resolverse a mano.
